% explainer-source-hash: f66797fe7a8c9d53245973007383b8fd1cc4dd9c9dfb65e80e4b5a7a95976308
% explainer-source-commit: a0ae5757644677d02ec29161881badb6a764a332
% explainer-generated: 2026-07-16
% Regenerate with /explain-all (see WIRING.md). Do not edit this stamp by hand:
% it is how the toolchain knows whether this document still matches the code.
\documentclass[11pt,a4paper]{article}
% preamble.tex - shared preamble for this project's explainer documents.
%
% This file is generated: it is copied in beside the documents on every build, so a
% clone of this repo can rebuild the PDFs, but local edits to it are overwritten.
%
% Usage from a document:
%   \documentclass[11pt,a4paper]{article}
%   % preamble.tex - shared preamble for this project's explainer documents.
%
% This file is generated: it is copied in beside the documents on every build, so a
% clone of this repo can rebuild the PDFs, but local edits to it are overwritten.
%
% Usage from a document:
%   \documentclass[11pt,a4paper]{article}
%   % preamble.tex - shared preamble for this project's explainer documents.
%
% This file is generated: it is copied in beside the documents on every build, so a
% clone of this repo can rebuild the PDFs, but local edits to it are overwritten.
%
% Usage from a document:
%   \documentclass[11pt,a4paper]{article}
%   \input{preamble}
%   \begin{document} ... \end{document}
%
% Build:  pdflatex -interaction=nonstopmode deep-dive.tex   (run it three times, so the
%         table of contents and the TikZ positioning settle)
%
% PACKAGE NOTES. Every package used here was checked present on the machine these
% documents were built on (Debian TeX Live 2023). If you rebuild elsewhere and a build
% fails, these are the likely gaps. Do not add a package without checking first:
%   kpsewhich <pkg>.sty
%
% Deliberately NOT used, because they were absent and the build dies without them:
%   algorithm2e, algorithm, algorithmic, algpseudocode  (texlive-science not installed)
%     -> the \begin{pseudocode} environment defined below stands in for them.
%   lmodern, charter, mathptmx, times, mathpazo, courier, newtx*, tgtermes
%     -> their .sty files RESOLVE but the font files themselves are absent, so the
%        build dies with "I can't find file `bchr8t'" or similar. Verified by
%        test-compiling each one. Only Computer Modern (the default) and helvet
%        actually render. kpsewhich reporting a font it cannot render is the trap:
%        do not "upgrade" the font stack without test-compiling it first.
%   amsmath, amssymb -> not loaded, so \text{} and \blacksquare are unavailable.
%
% A fuller font stack and algorithm2e would come from:
%   sudo apt install texlive-fonts-recommended texlive-science cm-super

\usepackage[T1]{fontenc}
\usepackage[utf8]{inputenc}
% Font: Computer Modern (default). See the note above before changing this.
\usepackage[a4paper,margin=2.4cm,headheight=14pt]{geometry}
% microtype WITHOUT font expansion. This is load-bearing, do not "restore" it:
% cm-super is not installed, so T1 Computer Modern resolves to BITMAP fonts, and
% pdfTeX's auto-expansion needs scalable ones. With expansion on, any document
% long enough to reach a second page dies with
%   "pdfTeX error (font expansion): auto expansion is only possible with scalable fonts"
% Protrusion (the bigger visual win) still works and stays on.
% `sudo apt install cm-super` would lift this; until then, leave expansion off.
\usepackage[protrusion=true,expansion=false]{microtype}
% Kill hyphen ligatures in TYPEWRITER text. This is load-bearing, not cosmetic. Under T1,
% \texttt{--verbose} silently renders as "-verbose", and \texttt{---} as "--", so a
% document can print a command-line flag as the WRONG flag and look perfectly fine doing
% it. Verified by rendering the page and reading the glyphs. lstlisting and \verb were
% always safe; this makes \code, \file and \texttt safe too, which is what prose uses.
% Ordinary prose is unaffected and still ligates, which is what you want there.
\DisableLigatures[-]{family=tt*}
\usepackage{graphicx}
\usepackage{xcolor}
\usepackage{booktabs}
\usepackage{enumitem}
\usepackage{fancyhdr}
\usepackage{titlesec}
\usepackage{tcolorbox}
\usepackage{listings}
\usepackage{float}
\usepackage{caption}
\usepackage{tikz}
\usepackage{forest}
\usepackage{pgfplots}
\pgfplotsset{compat=1.18}
\usetikzlibrary{arrows.meta,positioning,shapes.geometric,fit,backgrounds,calc,decorations.pathreplacing}
\tcbuselibrary{skins,breakable}
\usepackage[hidelinks,bookmarks=true,bookmarksnumbered=true]{hyperref}

% ===== palette ===========================================================
\definecolor{inkblue}{HTML}{1F3A5F}
\definecolor{accent}{HTML}{2E7DAF}
\definecolor{softgrey}{HTML}{F4F6F8}
\definecolor{linegrey}{HTML}{C8D0D8}
\definecolor{codebg}{HTML}{F7F7F4}
\definecolor{warnred}{HTML}{A33B3B}
\definecolor{okgreen}{HTML}{2E7D53}

% ===== headings ==========================================================
\titleformat{\section}{\Large\bfseries\color{inkblue}}{\thesection}{0.7em}{}
\titleformat{\subsection}{\large\bfseries\color{inkblue}}{\thesubsection}{0.6em}{}
\titleformat{\subsubsection}{\normalsize\bfseries\color{accent}}{\thesubsubsection}{0.5em}{}

\pagestyle{fancy}
\fancyhf{}
\fancyhead[L]{\small\color{inkblue}\nouppercase{\leftmark}}
\fancyhead[R]{\small\color{inkblue}\thepage}
\renewcommand{\headrulewidth}{0.4pt}
\renewcommand{\headrule}{\hbox to\headwidth{\color{linegrey}\leaders\hrule height \headrulewidth\hfill}}

% ===== code ==============================================================
\lstdefinestyle{house}{
  backgroundcolor=\color{codebg},
  basicstyle=\ttfamily\footnotesize,
  keywordstyle=\color{inkblue}\bfseries,
  commentstyle=\color{gray}\itshape,
  stringstyle=\color{okgreen},
  numbers=left, numberstyle=\tiny\color{gray}, numbersep=8pt,
  frame=single, rulecolor=\color{linegrey},
  breaklines=true, breakatwhitespace=false,
  showstringspaces=false, tabsize=2,
  captionpos=b, xleftmargin=14pt, framexleftmargin=14pt,
  columns=fullflexible, keepspaces=true,
  upquote=true,  % render ' and " as straight quotes, not curly
  % NON-ASCII INSIDE A LISTING IS A HARD BUILD FAILURE, not a cosmetic issue: listings
  % reads raw bytes, so inputenc dies on the first multi-byte character ("Invalid UTF-8
  % byte sequence") and the whole document fails to build. These mappings cover what
  % realistically turns up in source code: Latin-1 accents, Spanish punctuation, smart
  % quotes picked up from a word processor, the degree and section signs.
  % NOT COVERED, and it cannot be here: Urdu, Arabic and Balochi script. pdflatex has no
  % font for them. Describe or transliterate such text; never paste it into a listing.
  literate=%
    {á}{{\'a}}1 {é}{{\'e}}1 {í}{{\'i}}1 {ó}{{\'o}}1 {ú}{{\'u}}1
    {Á}{{\'A}}1 {É}{{\'E}}1 {Í}{{\'I}}1 {Ó}{{\'O}}1 {Ú}{{\'U}}1
    {à}{{\`a}}1 {è}{{\`e}}1 {ì}{{\`i}}1 {ò}{{\`o}}1 {ù}{{\`u}}1
    {ä}{{\"a}}1 {ë}{{\"e}}1 {ï}{{\"i}}1 {ö}{{\"o}}1 {ü}{{\"u}}1
    {Ä}{{\"A}}1 {Ö}{{\"O}}1 {Ü}{{\"U}}1 {ß}{{\ss}}1
    {ñ}{{\~n}}1 {Ñ}{{\~N}}1 {ç}{{\c c}}1 {Ç}{{\c C}}1
    {¿}{{\textquestiondown}}1 {¡}{{\textexclamdown}}1
    {°}{{\textdegree}}1 {§}{{\S}}1 {£}{{\pounds}}1 {€}{{EUR}}1
    {“}{{\textquotedblleft}}1 {”}{{\textquotedblright}}1
    {‘}{{\textquoteleft}}1 {’}{{\textquoteright}}1
    {–}{{-}}1 {—}{{-}}1 {…}{{...}}1 {→}{{$\rightarrow$}}1 {✓}{{x}}1
}
\lstset{style=house}

% ===== callout boxes =====================================================
% \begin{keyidea}{Title} ... \end{keyidea}   - a concept a beginner must grasp
\newtcolorbox{keyidea}[1]{
  enhanced, breakable, colback=softgrey, colframe=accent,
  boxrule=0.4pt, arc=2pt, left=8pt, right=8pt, top=6pt, bottom=6pt,
  title={\bfseries #1}, coltitle=white, colbacktitle=accent, fonttitle=\small
}

% \begin{jargon}{Term} definition \end{jargon}  - define a term at first use
\newtcolorbox{jargon}[1]{
  enhanced, breakable, colback=white, colframe=linegrey,
  boxrule=0.4pt, arc=2pt, left=8pt, right=8pt, top=6pt, bottom=6pt,
  title={\bfseries Jargon: #1}, coltitle=inkblue, colbacktitle=softgrey, fonttitle=\small
}

% \begin{honest}{Title} ... \end{honest}  - brutally honest finding (rule 18)
\newtcolorbox{honest}[1]{
  enhanced, breakable, colback=white, colframe=warnred,
  boxrule=0.6pt, arc=2pt, left=8pt, right=8pt, top=6pt, bottom=6pt,
  title={\bfseries #1}, coltitle=white, colbacktitle=warnred, fonttitle=\small
}

% \begin{pseudocode}{Caption} ... \end{pseudocode}
% Stands in for algorithm2e, which is NOT installed. Use \Step / \Indent inside.
\newtcolorbox{pseudocode}[1]{
  enhanced, breakable, colback=codebg, colframe=inkblue,
  boxrule=0.4pt, arc=2pt, left=10pt, right=8pt, top=6pt, bottom=6pt,
  fontupper=\ttfamily\small,
  title={\bfseries\rmfamily #1}, coltitle=white, colbacktitle=inkblue, fonttitle=\small
}
\newcommand{\Step}[1]{\par\noindent #1}
\newcommand{\Indent}[1]{\par\noindent\hspace*{2em}#1}

% ===== tikz house styles =================================================
\tikzset{
  box/.style   = {rectangle, rounded corners=2pt, draw=inkblue, fill=softgrey,
                  text=inkblue, align=center, inner sep=6pt, minimum height=9mm,
                  font=\small},
  extbox/.style= {rectangle, rounded corners=2pt, draw=linegrey, fill=white,
                  text=black!70, align=center, inner sep=6pt, minimum height=9mm,
                  font=\small, dashed},
  store/.style = {cylinder, shape border rotate=90, aspect=0.25, draw=accent,
                  fill=white, text=inkblue, align=center, inner sep=5pt,
                  font=\small},
  dec/.style   = {diamond, aspect=1.6, draw=accent, fill=white, text=inkblue,
                  align=center, inner sep=2pt, font=\scriptsize},
  flow/.style  = {-{Stealth[length=2.4mm]}, draw=inkblue, thick},
  dflow/.style = {-{Stealth[length=2.4mm]}, draw=accent, thick, dashed},
  % align=center is load-bearing: without it, a \\ inside an edge label throws
  % the misleading "perhaps a missing \item" error.
  lbl/.style   = {font=\scriptsize\itshape, text=black!65, fill=white,
                  inner sep=2pt, align=center}
}

% forest default for directory trees
\forestset{
  dirtree/.style={
    for tree={
      font=\ttfamily\small, grow'=0, child anchor=west, parent anchor=south,
      anchor=west, calign=first, inner sep=2pt,
      edge path={\noexpand\path [draw, \forestoption{edge}]
        (!u.south west) +(9pt,0) |- (.child anchor) \forestoption{edge label};},
      before typesetting nodes={if n=1{insert before={[,phantom]}}{}},
      fit=band, before computing xy={l=15pt}
    }
  }
}

% ===== misc ==============================================================
\captionsetup{font=small, labelfont={bf,color=inkblue}}
\setlist[itemize]{leftmargin=1.2em, itemsep=2pt, topsep=3pt}
\setlist[enumerate]{leftmargin=1.4em, itemsep=2pt, topsep=3pt}
\newcommand{\file}[1]{\texttt{\small #1}}
\newcommand{\code}[1]{\texttt{\small #1}}
\setcounter{tocdepth}{2}
\setcounter{secnumdepth}{3}

%   \begin{document} ... \end{document}
%
% Build:  pdflatex -interaction=nonstopmode deep-dive.tex   (run it three times, so the
%         table of contents and the TikZ positioning settle)
%
% PACKAGE NOTES. Every package used here was checked present on the machine these
% documents were built on (Debian TeX Live 2023). If you rebuild elsewhere and a build
% fails, these are the likely gaps. Do not add a package without checking first:
%   kpsewhich <pkg>.sty
%
% Deliberately NOT used, because they were absent and the build dies without them:
%   algorithm2e, algorithm, algorithmic, algpseudocode  (texlive-science not installed)
%     -> the \begin{pseudocode} environment defined below stands in for them.
%   lmodern, charter, mathptmx, times, mathpazo, courier, newtx*, tgtermes
%     -> their .sty files RESOLVE but the font files themselves are absent, so the
%        build dies with "I can't find file `bchr8t'" or similar. Verified by
%        test-compiling each one. Only Computer Modern (the default) and helvet
%        actually render. kpsewhich reporting a font it cannot render is the trap:
%        do not "upgrade" the font stack without test-compiling it first.
%   amsmath, amssymb -> not loaded, so \text{} and \blacksquare are unavailable.
%
% A fuller font stack and algorithm2e would come from:
%   sudo apt install texlive-fonts-recommended texlive-science cm-super

\usepackage[T1]{fontenc}
\usepackage[utf8]{inputenc}
% Font: Computer Modern (default). See the note above before changing this.
\usepackage[a4paper,margin=2.4cm,headheight=14pt]{geometry}
% microtype WITHOUT font expansion. This is load-bearing, do not "restore" it:
% cm-super is not installed, so T1 Computer Modern resolves to BITMAP fonts, and
% pdfTeX's auto-expansion needs scalable ones. With expansion on, any document
% long enough to reach a second page dies with
%   "pdfTeX error (font expansion): auto expansion is only possible with scalable fonts"
% Protrusion (the bigger visual win) still works and stays on.
% `sudo apt install cm-super` would lift this; until then, leave expansion off.
\usepackage[protrusion=true,expansion=false]{microtype}
% Kill hyphen ligatures in TYPEWRITER text. This is load-bearing, not cosmetic. Under T1,
% \texttt{--verbose} silently renders as "-verbose", and \texttt{---} as "--", so a
% document can print a command-line flag as the WRONG flag and look perfectly fine doing
% it. Verified by rendering the page and reading the glyphs. lstlisting and \verb were
% always safe; this makes \code, \file and \texttt safe too, which is what prose uses.
% Ordinary prose is unaffected and still ligates, which is what you want there.
\DisableLigatures[-]{family=tt*}
\usepackage{graphicx}
\usepackage{xcolor}
\usepackage{booktabs}
\usepackage{enumitem}
\usepackage{fancyhdr}
\usepackage{titlesec}
\usepackage{tcolorbox}
\usepackage{listings}
\usepackage{float}
\usepackage{caption}
\usepackage{tikz}
\usepackage{forest}
\usepackage{pgfplots}
\pgfplotsset{compat=1.18}
\usetikzlibrary{arrows.meta,positioning,shapes.geometric,fit,backgrounds,calc,decorations.pathreplacing}
\tcbuselibrary{skins,breakable}
\usepackage[hidelinks,bookmarks=true,bookmarksnumbered=true]{hyperref}

% ===== palette ===========================================================
\definecolor{inkblue}{HTML}{1F3A5F}
\definecolor{accent}{HTML}{2E7DAF}
\definecolor{softgrey}{HTML}{F4F6F8}
\definecolor{linegrey}{HTML}{C8D0D8}
\definecolor{codebg}{HTML}{F7F7F4}
\definecolor{warnred}{HTML}{A33B3B}
\definecolor{okgreen}{HTML}{2E7D53}

% ===== headings ==========================================================
\titleformat{\section}{\Large\bfseries\color{inkblue}}{\thesection}{0.7em}{}
\titleformat{\subsection}{\large\bfseries\color{inkblue}}{\thesubsection}{0.6em}{}
\titleformat{\subsubsection}{\normalsize\bfseries\color{accent}}{\thesubsubsection}{0.5em}{}

\pagestyle{fancy}
\fancyhf{}
\fancyhead[L]{\small\color{inkblue}\nouppercase{\leftmark}}
\fancyhead[R]{\small\color{inkblue}\thepage}
\renewcommand{\headrulewidth}{0.4pt}
\renewcommand{\headrule}{\hbox to\headwidth{\color{linegrey}\leaders\hrule height \headrulewidth\hfill}}

% ===== code ==============================================================
\lstdefinestyle{house}{
  backgroundcolor=\color{codebg},
  basicstyle=\ttfamily\footnotesize,
  keywordstyle=\color{inkblue}\bfseries,
  commentstyle=\color{gray}\itshape,
  stringstyle=\color{okgreen},
  numbers=left, numberstyle=\tiny\color{gray}, numbersep=8pt,
  frame=single, rulecolor=\color{linegrey},
  breaklines=true, breakatwhitespace=false,
  showstringspaces=false, tabsize=2,
  captionpos=b, xleftmargin=14pt, framexleftmargin=14pt,
  columns=fullflexible, keepspaces=true,
  upquote=true,  % render ' and " as straight quotes, not curly
  % NON-ASCII INSIDE A LISTING IS A HARD BUILD FAILURE, not a cosmetic issue: listings
  % reads raw bytes, so inputenc dies on the first multi-byte character ("Invalid UTF-8
  % byte sequence") and the whole document fails to build. These mappings cover what
  % realistically turns up in source code: Latin-1 accents, Spanish punctuation, smart
  % quotes picked up from a word processor, the degree and section signs.
  % NOT COVERED, and it cannot be here: Urdu, Arabic and Balochi script. pdflatex has no
  % font for them. Describe or transliterate such text; never paste it into a listing.
  literate=%
    {á}{{\'a}}1 {é}{{\'e}}1 {í}{{\'i}}1 {ó}{{\'o}}1 {ú}{{\'u}}1
    {Á}{{\'A}}1 {É}{{\'E}}1 {Í}{{\'I}}1 {Ó}{{\'O}}1 {Ú}{{\'U}}1
    {à}{{\`a}}1 {è}{{\`e}}1 {ì}{{\`i}}1 {ò}{{\`o}}1 {ù}{{\`u}}1
    {ä}{{\"a}}1 {ë}{{\"e}}1 {ï}{{\"i}}1 {ö}{{\"o}}1 {ü}{{\"u}}1
    {Ä}{{\"A}}1 {Ö}{{\"O}}1 {Ü}{{\"U}}1 {ß}{{\ss}}1
    {ñ}{{\~n}}1 {Ñ}{{\~N}}1 {ç}{{\c c}}1 {Ç}{{\c C}}1
    {¿}{{\textquestiondown}}1 {¡}{{\textexclamdown}}1
    {°}{{\textdegree}}1 {§}{{\S}}1 {£}{{\pounds}}1 {€}{{EUR}}1
    {“}{{\textquotedblleft}}1 {”}{{\textquotedblright}}1
    {‘}{{\textquoteleft}}1 {’}{{\textquoteright}}1
    {–}{{-}}1 {—}{{-}}1 {…}{{...}}1 {→}{{$\rightarrow$}}1 {✓}{{x}}1
}
\lstset{style=house}

% ===== callout boxes =====================================================
% \begin{keyidea}{Title} ... \end{keyidea}   - a concept a beginner must grasp
\newtcolorbox{keyidea}[1]{
  enhanced, breakable, colback=softgrey, colframe=accent,
  boxrule=0.4pt, arc=2pt, left=8pt, right=8pt, top=6pt, bottom=6pt,
  title={\bfseries #1}, coltitle=white, colbacktitle=accent, fonttitle=\small
}

% \begin{jargon}{Term} definition \end{jargon}  - define a term at first use
\newtcolorbox{jargon}[1]{
  enhanced, breakable, colback=white, colframe=linegrey,
  boxrule=0.4pt, arc=2pt, left=8pt, right=8pt, top=6pt, bottom=6pt,
  title={\bfseries Jargon: #1}, coltitle=inkblue, colbacktitle=softgrey, fonttitle=\small
}

% \begin{honest}{Title} ... \end{honest}  - brutally honest finding (rule 18)
\newtcolorbox{honest}[1]{
  enhanced, breakable, colback=white, colframe=warnred,
  boxrule=0.6pt, arc=2pt, left=8pt, right=8pt, top=6pt, bottom=6pt,
  title={\bfseries #1}, coltitle=white, colbacktitle=warnred, fonttitle=\small
}

% \begin{pseudocode}{Caption} ... \end{pseudocode}
% Stands in for algorithm2e, which is NOT installed. Use \Step / \Indent inside.
\newtcolorbox{pseudocode}[1]{
  enhanced, breakable, colback=codebg, colframe=inkblue,
  boxrule=0.4pt, arc=2pt, left=10pt, right=8pt, top=6pt, bottom=6pt,
  fontupper=\ttfamily\small,
  title={\bfseries\rmfamily #1}, coltitle=white, colbacktitle=inkblue, fonttitle=\small
}
\newcommand{\Step}[1]{\par\noindent #1}
\newcommand{\Indent}[1]{\par\noindent\hspace*{2em}#1}

% ===== tikz house styles =================================================
\tikzset{
  box/.style   = {rectangle, rounded corners=2pt, draw=inkblue, fill=softgrey,
                  text=inkblue, align=center, inner sep=6pt, minimum height=9mm,
                  font=\small},
  extbox/.style= {rectangle, rounded corners=2pt, draw=linegrey, fill=white,
                  text=black!70, align=center, inner sep=6pt, minimum height=9mm,
                  font=\small, dashed},
  store/.style = {cylinder, shape border rotate=90, aspect=0.25, draw=accent,
                  fill=white, text=inkblue, align=center, inner sep=5pt,
                  font=\small},
  dec/.style   = {diamond, aspect=1.6, draw=accent, fill=white, text=inkblue,
                  align=center, inner sep=2pt, font=\scriptsize},
  flow/.style  = {-{Stealth[length=2.4mm]}, draw=inkblue, thick},
  dflow/.style = {-{Stealth[length=2.4mm]}, draw=accent, thick, dashed},
  % align=center is load-bearing: without it, a \\ inside an edge label throws
  % the misleading "perhaps a missing \item" error.
  lbl/.style   = {font=\scriptsize\itshape, text=black!65, fill=white,
                  inner sep=2pt, align=center}
}

% forest default for directory trees
\forestset{
  dirtree/.style={
    for tree={
      font=\ttfamily\small, grow'=0, child anchor=west, parent anchor=south,
      anchor=west, calign=first, inner sep=2pt,
      edge path={\noexpand\path [draw, \forestoption{edge}]
        (!u.south west) +(9pt,0) |- (.child anchor) \forestoption{edge label};},
      before typesetting nodes={if n=1{insert before={[,phantom]}}{}},
      fit=band, before computing xy={l=15pt}
    }
  }
}

% ===== misc ==============================================================
\captionsetup{font=small, labelfont={bf,color=inkblue}}
\setlist[itemize]{leftmargin=1.2em, itemsep=2pt, topsep=3pt}
\setlist[enumerate]{leftmargin=1.4em, itemsep=2pt, topsep=3pt}
\newcommand{\file}[1]{\texttt{\small #1}}
\newcommand{\code}[1]{\texttt{\small #1}}
\setcounter{tocdepth}{2}
\setcounter{secnumdepth}{3}

%   \begin{document} ... \end{document}
%
% Build:  pdflatex -interaction=nonstopmode deep-dive.tex   (run it three times, so the
%         table of contents and the TikZ positioning settle)
%
% PACKAGE NOTES. Every package used here was checked present on the machine these
% documents were built on (Debian TeX Live 2023). If you rebuild elsewhere and a build
% fails, these are the likely gaps. Do not add a package without checking first:
%   kpsewhich <pkg>.sty
%
% Deliberately NOT used, because they were absent and the build dies without them:
%   algorithm2e, algorithm, algorithmic, algpseudocode  (texlive-science not installed)
%     -> the \begin{pseudocode} environment defined below stands in for them.
%   lmodern, charter, mathptmx, times, mathpazo, courier, newtx*, tgtermes
%     -> their .sty files RESOLVE but the font files themselves are absent, so the
%        build dies with "I can't find file `bchr8t'" or similar. Verified by
%        test-compiling each one. Only Computer Modern (the default) and helvet
%        actually render. kpsewhich reporting a font it cannot render is the trap:
%        do not "upgrade" the font stack without test-compiling it first.
%   amsmath, amssymb -> not loaded, so \text{} and \blacksquare are unavailable.
%
% A fuller font stack and algorithm2e would come from:
%   sudo apt install texlive-fonts-recommended texlive-science cm-super

\usepackage[T1]{fontenc}
\usepackage[utf8]{inputenc}
% Font: Computer Modern (default). See the note above before changing this.
\usepackage[a4paper,margin=2.4cm,headheight=14pt]{geometry}
% microtype WITHOUT font expansion. This is load-bearing, do not "restore" it:
% cm-super is not installed, so T1 Computer Modern resolves to BITMAP fonts, and
% pdfTeX's auto-expansion needs scalable ones. With expansion on, any document
% long enough to reach a second page dies with
%   "pdfTeX error (font expansion): auto expansion is only possible with scalable fonts"
% Protrusion (the bigger visual win) still works and stays on.
% `sudo apt install cm-super` would lift this; until then, leave expansion off.
\usepackage[protrusion=true,expansion=false]{microtype}
% Kill hyphen ligatures in TYPEWRITER text. This is load-bearing, not cosmetic. Under T1,
% \texttt{--verbose} silently renders as "-verbose", and \texttt{---} as "--", so a
% document can print a command-line flag as the WRONG flag and look perfectly fine doing
% it. Verified by rendering the page and reading the glyphs. lstlisting and \verb were
% always safe; this makes \code, \file and \texttt safe too, which is what prose uses.
% Ordinary prose is unaffected and still ligates, which is what you want there.
\DisableLigatures[-]{family=tt*}
\usepackage{graphicx}
\usepackage{xcolor}
\usepackage{booktabs}
\usepackage{enumitem}
\usepackage{fancyhdr}
\usepackage{titlesec}
\usepackage{tcolorbox}
\usepackage{listings}
\usepackage{float}
\usepackage{caption}
\usepackage{tikz}
\usepackage{forest}
\usepackage{pgfplots}
\pgfplotsset{compat=1.18}
\usetikzlibrary{arrows.meta,positioning,shapes.geometric,fit,backgrounds,calc,decorations.pathreplacing}
\tcbuselibrary{skins,breakable}
\usepackage[hidelinks,bookmarks=true,bookmarksnumbered=true]{hyperref}

% ===== palette ===========================================================
\definecolor{inkblue}{HTML}{1F3A5F}
\definecolor{accent}{HTML}{2E7DAF}
\definecolor{softgrey}{HTML}{F4F6F8}
\definecolor{linegrey}{HTML}{C8D0D8}
\definecolor{codebg}{HTML}{F7F7F4}
\definecolor{warnred}{HTML}{A33B3B}
\definecolor{okgreen}{HTML}{2E7D53}

% ===== headings ==========================================================
\titleformat{\section}{\Large\bfseries\color{inkblue}}{\thesection}{0.7em}{}
\titleformat{\subsection}{\large\bfseries\color{inkblue}}{\thesubsection}{0.6em}{}
\titleformat{\subsubsection}{\normalsize\bfseries\color{accent}}{\thesubsubsection}{0.5em}{}

\pagestyle{fancy}
\fancyhf{}
\fancyhead[L]{\small\color{inkblue}\nouppercase{\leftmark}}
\fancyhead[R]{\small\color{inkblue}\thepage}
\renewcommand{\headrulewidth}{0.4pt}
\renewcommand{\headrule}{\hbox to\headwidth{\color{linegrey}\leaders\hrule height \headrulewidth\hfill}}

% ===== code ==============================================================
\lstdefinestyle{house}{
  backgroundcolor=\color{codebg},
  basicstyle=\ttfamily\footnotesize,
  keywordstyle=\color{inkblue}\bfseries,
  commentstyle=\color{gray}\itshape,
  stringstyle=\color{okgreen},
  numbers=left, numberstyle=\tiny\color{gray}, numbersep=8pt,
  frame=single, rulecolor=\color{linegrey},
  breaklines=true, breakatwhitespace=false,
  showstringspaces=false, tabsize=2,
  captionpos=b, xleftmargin=14pt, framexleftmargin=14pt,
  columns=fullflexible, keepspaces=true,
  upquote=true,  % render ' and " as straight quotes, not curly
  % NON-ASCII INSIDE A LISTING IS A HARD BUILD FAILURE, not a cosmetic issue: listings
  % reads raw bytes, so inputenc dies on the first multi-byte character ("Invalid UTF-8
  % byte sequence") and the whole document fails to build. These mappings cover what
  % realistically turns up in source code: Latin-1 accents, Spanish punctuation, smart
  % quotes picked up from a word processor, the degree and section signs.
  % NOT COVERED, and it cannot be here: Urdu, Arabic and Balochi script. pdflatex has no
  % font for them. Describe or transliterate such text; never paste it into a listing.
  literate=%
    {á}{{\'a}}1 {é}{{\'e}}1 {í}{{\'i}}1 {ó}{{\'o}}1 {ú}{{\'u}}1
    {Á}{{\'A}}1 {É}{{\'E}}1 {Í}{{\'I}}1 {Ó}{{\'O}}1 {Ú}{{\'U}}1
    {à}{{\`a}}1 {è}{{\`e}}1 {ì}{{\`i}}1 {ò}{{\`o}}1 {ù}{{\`u}}1
    {ä}{{\"a}}1 {ë}{{\"e}}1 {ï}{{\"i}}1 {ö}{{\"o}}1 {ü}{{\"u}}1
    {Ä}{{\"A}}1 {Ö}{{\"O}}1 {Ü}{{\"U}}1 {ß}{{\ss}}1
    {ñ}{{\~n}}1 {Ñ}{{\~N}}1 {ç}{{\c c}}1 {Ç}{{\c C}}1
    {¿}{{\textquestiondown}}1 {¡}{{\textexclamdown}}1
    {°}{{\textdegree}}1 {§}{{\S}}1 {£}{{\pounds}}1 {€}{{EUR}}1
    {“}{{\textquotedblleft}}1 {”}{{\textquotedblright}}1
    {‘}{{\textquoteleft}}1 {’}{{\textquoteright}}1
    {–}{{-}}1 {—}{{-}}1 {…}{{...}}1 {→}{{$\rightarrow$}}1 {✓}{{x}}1
}
\lstset{style=house}

% ===== callout boxes =====================================================
% \begin{keyidea}{Title} ... \end{keyidea}   - a concept a beginner must grasp
\newtcolorbox{keyidea}[1]{
  enhanced, breakable, colback=softgrey, colframe=accent,
  boxrule=0.4pt, arc=2pt, left=8pt, right=8pt, top=6pt, bottom=6pt,
  title={\bfseries #1}, coltitle=white, colbacktitle=accent, fonttitle=\small
}

% \begin{jargon}{Term} definition \end{jargon}  - define a term at first use
\newtcolorbox{jargon}[1]{
  enhanced, breakable, colback=white, colframe=linegrey,
  boxrule=0.4pt, arc=2pt, left=8pt, right=8pt, top=6pt, bottom=6pt,
  title={\bfseries Jargon: #1}, coltitle=inkblue, colbacktitle=softgrey, fonttitle=\small
}

% \begin{honest}{Title} ... \end{honest}  - brutally honest finding (rule 18)
\newtcolorbox{honest}[1]{
  enhanced, breakable, colback=white, colframe=warnred,
  boxrule=0.6pt, arc=2pt, left=8pt, right=8pt, top=6pt, bottom=6pt,
  title={\bfseries #1}, coltitle=white, colbacktitle=warnred, fonttitle=\small
}

% \begin{pseudocode}{Caption} ... \end{pseudocode}
% Stands in for algorithm2e, which is NOT installed. Use \Step / \Indent inside.
\newtcolorbox{pseudocode}[1]{
  enhanced, breakable, colback=codebg, colframe=inkblue,
  boxrule=0.4pt, arc=2pt, left=10pt, right=8pt, top=6pt, bottom=6pt,
  fontupper=\ttfamily\small,
  title={\bfseries\rmfamily #1}, coltitle=white, colbacktitle=inkblue, fonttitle=\small
}
\newcommand{\Step}[1]{\par\noindent #1}
\newcommand{\Indent}[1]{\par\noindent\hspace*{2em}#1}

% ===== tikz house styles =================================================
\tikzset{
  box/.style   = {rectangle, rounded corners=2pt, draw=inkblue, fill=softgrey,
                  text=inkblue, align=center, inner sep=6pt, minimum height=9mm,
                  font=\small},
  extbox/.style= {rectangle, rounded corners=2pt, draw=linegrey, fill=white,
                  text=black!70, align=center, inner sep=6pt, minimum height=9mm,
                  font=\small, dashed},
  store/.style = {cylinder, shape border rotate=90, aspect=0.25, draw=accent,
                  fill=white, text=inkblue, align=center, inner sep=5pt,
                  font=\small},
  dec/.style   = {diamond, aspect=1.6, draw=accent, fill=white, text=inkblue,
                  align=center, inner sep=2pt, font=\scriptsize},
  flow/.style  = {-{Stealth[length=2.4mm]}, draw=inkblue, thick},
  dflow/.style = {-{Stealth[length=2.4mm]}, draw=accent, thick, dashed},
  % align=center is load-bearing: without it, a \\ inside an edge label throws
  % the misleading "perhaps a missing \item" error.
  lbl/.style   = {font=\scriptsize\itshape, text=black!65, fill=white,
                  inner sep=2pt, align=center}
}

% forest default for directory trees
\forestset{
  dirtree/.style={
    for tree={
      font=\ttfamily\small, grow'=0, child anchor=west, parent anchor=south,
      anchor=west, calign=first, inner sep=2pt,
      edge path={\noexpand\path [draw, \forestoption{edge}]
        (!u.south west) +(9pt,0) |- (.child anchor) \forestoption{edge label};},
      before typesetting nodes={if n=1{insert before={[,phantom]}}{}},
      fit=band, before computing xy={l=15pt}
    }
  }
}

% ===== misc ==============================================================
\captionsetup{font=small, labelfont={bf,color=inkblue}}
\setlist[itemize]{leftmargin=1.2em, itemsep=2pt, topsep=3pt}
\setlist[enumerate]{leftmargin=1.4em, itemsep=2pt, topsep=3pt}
\newcommand{\file}[1]{\texttt{\small #1}}
\newcommand{\code}[1]{\texttt{\small #1}}
\setcounter{tocdepth}{2}
\setcounter{secnumdepth}{3}


\begin{document}

\begin{center}
{\LARGE\bfseries\color{inkblue} ZaraiLink: Demo\par}
\vspace{4pt}
{\large\color{accent} Brief\par}
\vspace{10pt}
{\small A short, accurate account of what this repository is, what it proves, and where it is weak.\par}
\end{center}

\vspace{6pt}

\begin{honest}{The two things that must be said first}
\textbf{The data is fabricated.} Every company, contact, phone number, price and shipment
figure in this repository is invented. No real trade data appears anywhere, and no code
here makes a network request. This is an imitation of another application's screens, not
that application.

\medskip
\textbf{ZaraiLink is a team project.} Three people built it: Salman Adnan led it, and Umar
Kashif and Fahad Nadeem wrote real, substantial parts. The real system is private and none
of its source is in this repository. This demo is Salman's own work; the product it depicts
is not his alone.
\end{honest}

% =====================================================================
\section{The problem this repository solves}

ZaraiLink is an agricultural trade-intelligence platform: a Django and DRF backend, a React
frontend, a natural-language search pipeline, hybrid retrieval, and a learning-to-rank
model over Pakistani customs data. It is private, and it cannot be made public by one
author acting alone.

That creates a real bind. A hiring panel cannot be handed the source, and a screenshot
proves nothing. The answer taken here is to rebuild the \emph{screens} against
\emph{invented} data. The visual design language and the interaction patterns are things
any browser loading the real app already downloads and renders. The private parts stay out.

\begin{figure}[H]
\centering
\resizebox{0.94\textwidth}{!}{
\begin{tikzpicture}
  \node[extbox, minimum width=48mm, minimum height=26mm, align=center] (real)
    {\textbf{ZaraiLink}\\[3pt]\footnotesize private, 3 authors\\\footnotesize Django + DRF, React\\\footnotesize NLU, retrieval, ranking\\\footnotesize real customs data};
  \node[box, minimum width=48mm, minimum height=26mm, right=44mm of real, align=center] (demo)
    {\textbf{zarailink-demo}\\[3pt]\footnotesize public, Salman alone\\\footnotesize HTML, CSS, vanilla JS\\\footnotesize no backend, no build\\\footnotesize fabricated data};
  \draw[flow] ($(real.east)+(0,7mm)$) -- node[lbl, above, align=center] {screens, tokens,\\interactions} ($(demo.west)+(0,7mm)$);
  \draw[dflow, draw=warnred] ($(real.east)+(0,-5mm)$) -- node[lbl, below, align=center, text=warnred] {logic, models, data\\\textbf{not copied}} ($(demo.west)+(0,-5mm)$);
\end{tikzpicture}}
\caption{The demo copies the surface and deliberately copies none of the private logic.}
\end{figure}

% =====================================================================
\section{What it is, technically}

Five files, about 1,100 lines, no framework and no build step. What sits in the repository
is exactly what the browser runs; you can open \file{index.html} from a file path, offline,
and it works.

\begin{center}
\begin{tabular}{@{}llr@{}}
\toprule
\textbf{File} & \textbf{Role} & \textbf{Lines} \\
\midrule
\file{app.js}       & Routing, filtering, unlocking, all rendering & 399 \\
\file{styles.css}   & Colour tokens and component layout & 369 \\
\file{index.html}   & Six screens as static markup & 271 \\
\file{demo-data.js} & The fabricated dataset & 70 \\
\file{icons.js}     & A 20-entry SVG icon lookup & 34 \\
\bottomrule
\end{tabular}
\end{center}

Six screens live in \file{index.html} at once. CSS hides all but the one carrying an
\code{is-active} class, and navigation is just moving that class. All behaviour hangs off a
\textbf{single} click listener on \code{<body>} that inspects what was clicked.

\begin{figure}[H]
\centering
\resizebox{0.96\textwidth}{!}{
\begin{tikzpicture}[node distance=11mm]
  \node[box] (dash) {\textbf{Dashboard}\\\footnotesize search home};
  \node[box, right=22mm of dash] (results) {\textbf{Results}\\\footnotesize Data Dashboard};
  \node[box, below=of dash] (dir) {\textbf{Directory}};
  \node[box, right=22mm of dir] (profile) {\textbf{Profile}\\\footnotesize contact unlock};
  \node[box, below=of dir] (watch) {\textbf{Watchlist}};
  \node[box, right=22mm of watch] (sub) {\textbf{Subscription}};
  \draw[flow] (dash) -- node[lbl, above] {search} (results);
  \draw[flow] (dash) -- node[lbl, left] {nav} (dir);
  \draw[flow] (dir) -- node[lbl, above] {open} (profile);
  \draw[flow] (results) -- node[lbl, right] {card} (profile);
  \draw[flow] (dir) -- node[lbl, left] {nav} (watch);
  \draw[flow] (watch) -- node[lbl, above] {nav} (sub);
\end{tikzpicture}}
\caption{Six screens, one page. No URL ever changes.}
\end{figure}

\begin{keyidea}{Why event delegation is the right call here}
Every grid on these screens is drawn by writing an HTML string into \code{innerHTML}, which
destroys and recreates the elements. Any listener attached directly to a button would be
thrown away on every render. Listening once on \code{<body>} and using
\code{e.target.closest(...)} to find what was hit means re-rendering is free and breaks
nothing. It is the best structural decision in the codebase.
\end{keyidea}

% =====================================================================
\section{What actually works}

\begin{itemize}
\item \textbf{Per-contact unlocking.} The real application unlocks contacts one at a time,
  not per company, and the demo reproduces that faithfully. Each contact carries its own
  \code{unlocked} flag; a confirm modal stashes the pending target, re-checks the token
  balance at confirm time rather than trusting the earlier check, spends one token, flips
  exactly one contact, and re-renders from the model. This is the demo's centrepiece and it
  is correct.
\item \textbf{Masking.} Locked phone and email values are rendered mangled (first character,
  then bullets, domain preserved) and blurred with CSS, then flip to their real fabricated
  values only for the specific contact that was paid for.
\item \textbf{Seniority badges.} A short ordered chain of regular expressions turns a job
  title into one of five badges. It is genuinely reading the data, not hard-coded.
\item \textbf{The directory search box.} Filters company names case-insensitively on every
  keystroke, with an automatic empty state.
\item \textbf{The honesty scaffolding.} A yellow banner on every screen, a footer
  disclaimer, \code{noindex} in the head, all domains under the reserved
  \code{.example.com}, and the word ``fabricated'' printed into the results screen itself.
\item \textbf{No build step, no network, no secrets.} A scan for keys, tokens, passwords and
  connection strings returned nothing, as expected for a codebase that talks to nobody.
\end{itemize}

% =====================================================================
\section{Where it is weaker than it looks}

This section is the point of the document. Everything below was verified by reading the
source and, where it involved logic, by executing it.

\subsection{Two real bugs}

\begin{honest}{The category paywall advertises 500 tokens and charges 5}
The button reads ``Unlock for 500 Tokens''. The code guards on \code{tokenBalance < 5} and
then deducts 5. The user starts with 7, so the button \emph{works}, which is precisely why
the bug survives: nothing looks broken. Against the plan grants (Starter gives 20 tokens
per month) a 500-token unlock would be absurd, so the label is almost certainly the typo.
The cost 5 is also a magic number written twice, with the advertised price a third,
unrelated literal.
\end{honest}

\begin{honest}{Un-starring on the Watchlist screen does nothing visible}
The star handler unconditionally calls \code{renderDirectory()} and never
\code{renderWatchlist()}. Both grids share the same card template, so the Watchlist has star
buttons. Click one there and the model updates correctly, but the re-render lands on a
hidden screen; the grid in front of you never changes. Navigate away and back and the change
appears, which is the signature of this class of bug.

This is the most likely defect for an interviewer to hit by accident, because starring
things and then going to the watchlist to unstar one is the obvious way to explore the
feature.
\end{honest}

\subsection{The demo is thinner than its screens imply}

\begin{itemize}
\item \textbf{Search does not filter.} The query is used only as a heading. Search ``PVC
  Resin'' and you get the same four dextrose and urea companies as searching ``dextrose''.
  Any numeric query hard-codes the heading ``HS 1702.3000: Dextrose'' regardless of what was
  typed.
\item \textbf{Four trade-direction pills, two behaviours.} ``Foreign Suppliers'' and
  ``Foreign Buyers'' have \emph{identical} predicates (\code{country !== "Pakistan"}), as do
  the two Pakistani pills. Only the country axis does anything; the buyer/supplier axis is
  decoration. Verified by running them: 3, 3, 1, 1 results, the same 3 and the same 1.
\item \textbf{The Trade Directory has one mode wearing two hats.} Switching to ``Find
  Buyers'' rewrites the heading and the lede and shows the same four companies. There is no
  buyer/supplier distinction in the data. The empty state still says ``No suppliers found.''
\item \textbf{The shipment chip contradicts the data.} It shows
  \code{COMPANIES.length * 14}, which is 56. The dataset's own \code{shipments} fields sum
  to 125.
\item \textbf{Buying a plan is an unconditional faucet.} No payment, no current-plan state,
  no limit. Click Enterprise ten times for 3,000 tokens.
\item \textbf{The world is small.} Four companies, five contacts, one of them already free.
  With a starting balance of 7 a user can unlock everything and still have 3 left, so the
  scarcity the screens dramatise never bites.
\end{itemize}

\subsection{Dead controls, the highest-cost defect here}

Six visible things respond to the mouse and do nothing:

\begin{center}
\begin{tabular}{@{}p{5.6cm}p{9.4cm}@{}}
\toprule
\textbf{Control} & \textbf{Reality} \\
\midrule
Import/Export toggle & \code{currentScope} is written on every click and read \textbf{nowhere} \\
Theme toggle (moon) & Styled, \code{aria-label}led, no listener, and no dark theme exists \\
``I want to buy'' / ``I want to sell'' & No \code{data-view}, no handler. The two pills beside them \emph{do} navigate, which makes the dead pair harder to spot \\
User menu & \code{cursor: pointer} and a hover state advertise a dropdown that does not exist \\
``Refine by Product'' checkboxes & Real, correctly computed counts; no listener. Convincing precisely \emph{because} the counts are right \\
\bottomrule
\end{tabular}
\end{center}

For a demo whose entire purpose is to be clicked by a stranger, each dead control is an
invitation to find the seam. A demo with fewer controls that all work is strictly stronger
than one with more controls where half are scenery.

\subsection{One feature the README claims that the code cannot do}

\begin{honest}{``Product Search'' mode is unreachable}
\code{detectMode(query, scopeSelected)} returns \code{"product"} only when
\code{scopeSelected} is falsy, but \textbf{both} call sites pass the literal \code{true}. So
a query starting with a digit gives \code{"hscode"}, any other non-empty query gives
\code{"ai"}, and the empty case is caught by an early return that clears the pill entirely.
The green ``Product Search'' pill can never render. Confirmed by executing the function.

The README states the pill ``flips between HS Code Mode, Product Search, and AI Query
Mode''. Two of the three are real. Passing \code{currentScope} instead of \code{true} would
not fix it either, since it is initialised to \code{"import"} and is always truthy; the
intended design presumably starts the scope unselected.
\end{honest}

\subsection{Where the README and the code disagree}

\begin{honest}{The README describes a font strategy the code abandoned}
The README says the demo ``falls back to the closest system font instead of fetching''
Inter externally. It does not. It ships \file{assets/fonts/inter-variable.woff2} and uses
real Inter. The git history explains the drift: the commit that self-hosted the font came
after the README paragraph was written, and the paragraph was never updated.
\file{styles.css}'s own header comment is correct. \textbf{The code wins; the README is
stale.}

This matters more than a normal doc bug. A project sold on honesty about what is real
cannot afford a README that overstates its own code, however innocently it drifted.
\end{honest}

\subsection{Dead data}

\code{HS\_TREE} (a real fragment of the product's model: code, name, category) is declared
and referenced nowhere; the Dashboard hard-codes the example \code{1702.3000} in HTML
instead of reading it. \code{company.website} and \code{company.last\_active} are populated
on every company and rendered never. Ten of the twenty icons are never called, and several
are duplicated as inline markup in \file{index.html}, so one icon set is maintained in two
places.

\subsection{Security}

Nearly everything renders through \code{innerHTML}. \code{escapeHtml} handles \code{\&},
\code{<} and \code{>} but \textbf{not quotes}, and several values land inside HTML
attributes where a quote would break out.

This is \textbf{not exploitable}, and the reason is worth stating precisely: every string in
the application is a hard-coded literal. There is no user-supplied data, no network input,
no URL parameter, no storage, and the two text inputs are read but never rendered back as
HTML. The attack surface is empty. It is not a vulnerability; it is a habit that would
become one the instant this code met real data, which is exactly what the real ZaraiLink
does.

Likewise the paywall blur is pure CSS, so the values sit in the DOM for anyone who opens
the developer tools. Harmless here because everything behind it is invented, but a real
paywall must withhold data at the server and never send it.

% =====================================================================
\section{What is claimed but cannot be checked here}

\begin{honest}{Three claims rest outside this repository}
\begin{enumerate}
\item \textbf{``64 of 97 commits''} (README). The private repository's git history is not
  here. This is the owner's to stand behind; an interviewer who asks should be pointed at
  the private repository.
\item \textbf{That the screens match the real components.} The README and
  \file{styles.css} cite specific real files (\code{Dashboard.js},
  \code{DataDashboard.js}, \code{CompanyProfile.js}) and state the colours were copied 1:1
  from the real \code{index.css}. None of those files are in this repository, so the fidelity
  claim is \textbf{inferred, not verified}.
\item \textbf{That masking mirrors the real \code{maskValue()}.} Same reason. Inferred.
\end{enumerate}
Separately, the font file is 48 KB, which is small for a full variable Inter, so it is
\emph{probably} a subset. That is \textbf{inferred from file size} and was not confirmed by
inspecting the font tables.
\end{honest}

Note also that this demo is evidence of UI reconstruction skill and product judgment. It is
\textbf{not} evidence that the real backend works, because none of it is here.

% =====================================================================
\section{Licence and deployment}

\file{netlify.toml} is two lines (\code{publish = "."}) with no build command, because
there is nothing to build. \file{index.html} carries \code{noindex}, correctly: a page full
of invented companies and invented contact details should not appear in a search for those
names.

The licence is \textbf{PolyForm Strict 1.0.0}, which is source-available, \emph{not} open
source: it permits reading and nothing else. For portfolio work that is coherent, and
especially so here, where the underlying product is a private team project. It should not be
described as open source. The copyright line reads ``Copyright 2026 Salman Adnan'' on a
repository whose subject matter belongs to three people; the demo's code is Salman's alone,
so that holds, but the distinction between ``this demo's code'' and ``ZaraiLink'' is doing
real work in that line.

% =====================================================================
\section{The highest-value hour of work}

Not a new feature.

\begin{enumerate}
\item \textbf{Fix the 500/5 token label and the watchlist re-render.} Both are one-line
  fixes; both are things a clicking interviewer will find.
\item \textbf{Correct the README's font paragraph and soften the three-modes claim.}
\item \textbf{Wire up or delete the dead controls.} Deleting them costs nothing and removes
  six chances for the demo to look broken rather than scripted.
\end{enumerate}

The README's ``What I'd do differently'' proposes adding the Trade Intelligence cluster
(link prediction, node2vec, the GNN work). That is the right \emph{ambition} and the wrong
\emph{next step}: adding a sixth screen to a demo with six dead controls widens the seam
rather than closing it.

% =====================================================================
\section{Glossary}

\begin{description}[leftmargin=0pt, style=nextline]
\item[Backend / frontend] The frontend runs in the browser and draws what you see; the
  backend runs on a server and holds the data. This demo has only a frontend, and it asks
  nobody anything.
\item[Build step] A program that translates source the browser cannot run (JSX,
  TypeScript, Sass) into plain HTML, CSS and JS. This demo has none, which is why it still
  works when opened from a file path with no toolchain installed.
\item[DOM] The browser's live, in-memory tree of the page. JavaScript changes what you see
  by changing this tree.
\item[Event delegation] Listening for clicks once at the top of the page and working out
  what was hit, instead of attaching a listener to every button. Relies on \emph{bubbling},
  the browser's habit of firing an event on an element and then on each of its ancestors.
\item[HS code] The Harmonized System number for a traded good. Customs records are indexed
  by it. \code{1702.3000} is dextrose.
\item[Regular expression] A small pattern language for matching text. \code{/\^{}\textbackslash d/}
  asks ``does this start with a digit?''
\item[SVG] Pictures drawn with maths rather than pixels, so they stay sharp at any size.
  \code{currentColor} means ``whatever the CSS text colour is here'', which lets one icon
  definition work on a dark navbar and a white card.
\item[Variable font] One font file that renders at any weight on a continuous axis, rather
  than one file per weight.
\item[XSS] Cross-site scripting: dropping untrusted text into a page as HTML so that
  \code{<script>} in that text actually executes. Prevented by escaping.
\item[Design tokens] Named, reusable values (colours, radii, widths) that encode a design
  system's decisions. Here they are CSS custom properties in a \code{:root} block.
\item[PolyForm Strict] A source-available licence granting permission to read the source
  and nothing else. Not open source.
\end{description}

\end{document}
